\scanpage{80}\mspage{76}
La structure $\mathrel{|\unsure{\circ|0}|}$ sur un Spatios permet de récupérer la structure $\leq,C$ du spatios, \unsure{grace|grâce} à la proposition suivante :

% slanted
\margin{Décrire la relation d'ordre \unsure{can.} d'un disjoncteur : $x\leq y\Longleftrightarrow\forall z,\ z\mathrel{|\circ|}y\Longrightarrow z\mathrel{|\circ|}x$}

\uline{Proposition 1} Soit $(\deleted{\unsure{S}}\unsure{\mathbb{\Sigma}|\Sigma},\leq,C)$ un spatios, $\mathrel{|\circ|}$ la structure de disjonction associée. Alors l'application canonique
\[
\sigma:\Sigma\longrightarrow \mathbb{\Sigma}_{\Sigma}
\]
\uline{est} \deleted{\unsure{bijec|bij}} \uline{un isomorphisme de spatios}.

L'\deleted{\unsure{appl}}isomorphisme inverse est donné par
\[
\sigma^{-1}:S\longmapsto\operatorname{Sup}S.
\]

% slanted
\margin{\uline{NB} La relation d'ordre est exprimée en termes de $\mathrel{|\circ|}$ comme ci-dessus. D'autre part, $Cx=\underset{\substack{x'\in\Sigma\\\text{t.q. }x'\mathrel{|\circ|}x}}{\operatorname{Sup}}x'$}

\uline{Scholie} Pour un ensemble donné $\Sigma$, il y a corr. 1-1 \unsure{dentre|entre} les deux ensembles suivants :

(i) L'\unsure{ens.|Ens.} \inserted{$\mathscr{S}$} des structures de spatios \inserted{$(\leq,C)$} sur $\deleted{\unsure{E}}\Sigma$.

(ii) L'\unsure{ens|en} \inserted{$\mathscr{D}$} des structures de disjonction \inserted{$\mathrel{|\circ|}$} sur $\Sigma$, telles que l'application canonique
\[
\Sigma\longrightarrow\mathbb{\Sigma}_{\underline{\Sigma}}\qquad
(\text{où }\underline{\Sigma}\unsure{=|\simeq}(\Sigma,\mathrel{|\circ|}))
\tag*{\makebox[0pt][r]{\makebox[\displaywidth][l]{($\unsure{\ast|\times}$)}}}
\]
soit bijective.

\deleted{\unsure{A l'appli|A l'appl}} La bijection $\mathscr{S}\overset{\sim}{\longrightarrow}\mathscr{D}$ \deleted{\unsure{est donnée}} \inserted{\unsure{associe|associé}, à} \deleted{par la} une structure de spatios, la structure de disj. associée. La bijection inverse associe à une structure de disjonction $\in\mathscr{D}$, la structure de spatios sur $\Sigma$ \unsure{déduit|déduite} par tr. de str. de la bijection ($\ast$).

% archive sheet number: 79
